\documentclass[11pt, a4paper]{article}
\usepackage[utf8]{inputenc}
\usepackage{geometry}
\usepackage{amsmath, amssymb}
\usepackage{graphicx}
\usepackage{hyperref}
\usepackage{titlesec}
\usepackage{setspace}
\usepackage{mathpazo}
\usepackage{booktabs}

% Geometry setup
\geometry{top=2.5cm, bottom=2.5cm, left=2.5cm, right=2.5cm}

% Formatting
\titleformat{\section}{\large\bfseries}{\thesection}{1em}{}
\onehalfspacing

% Colors for hyperlinks
\hypersetup{
    colorlinks=true,
    linkcolor=blue,
    filecolor=magenta,
    urlcolor=cyan,
    citecolor=blue,
}

\title{\textbf{From Kinematics to Hazard: \\ A Unified Physics-Informed Inversion of the Iranian Lithosphere}}
\author{\textbf{Research Proposal (Version 4.0 Draft)}}
\date{\today}

% ============================================================================
% Revision history for v4
% ============================================================================
% v4 is being built incrementally from v3, one amendment at a time. Each
% amendment corresponds to a leaf in the phase plan documented in
% audits/2026-06-03-foundation.md.
%
% Applied amendments (committed):
%   P.1  Locked the Monte Carlo estimator form of L_seis (Eq. 7 -> Eq. 7+8).
%   P.2  z-coordinate normalization convention: [-1, 1]^2 x [-1, 0] with
%        surface at z=0 and max depth at z=-1 (new Section 3.3).
%
% Pending amendments (queued for review):
%   P.3  UTM zone handling (single vs. multi)
%   P.4  Magnitude unification rule (Mw priority + Scordilis conversions)
%   P.5  Magnitude of completeness (MAXC + GFT cross-check)
%   P.6  Declustering parameters (Zaliapin-Ben-Zion)
%   P.7  Constitutive law: viscous-only (remove elastic)
%   P.8  Network hyperparameter formalization
%   P.9  Data sufficiency footnote (209 GPS vs. 2000)
%   P.10 Random seed policy
% ============================================================================

\begin{document}

\maketitle

\begin{abstract}
Mapping the crustal stress field and forecasting earthquakes have traditionally been treated as separate problems: geodesists invert GPS velocities for strain, while seismologists compute probabilistic hazard from earthquake catalogs. We introduce a \textbf{Unified Physics-Informed Neural Network (PINN)} that couples both within a single trainable framework. Our network acts as a mesh-free solver for the 3D momentum balance equation, constrained by heterogeneous material properties derived from P-wave tomography. Unlike previous PINN stress inversions \cite{Masson2023}, which are limited to 2D homogeneous elasticity, our model (i) operates in three dimensions with spatially varying rheology, (ii) uses Fourier feature embeddings to resolve high-frequency stress gradients, and (iii) introduces a \textbf{Coulomb--Dieterich seismicity coupling} that anchors the absolute stress magnitude to the observed earthquake catalog via rate-and-state friction physics \cite{Dieterich1994}. This coupling resolves the fundamental magnitude ambiguity of kinematic-only inversions. Applied to Central Iran, the framework produces the first continuous, physics-constrained 3D map of Coulomb seismogenic potential for the Iranian Plateau.
\end{abstract}

\section{The Core Problem: Magnitude Ambiguity in Stress Inversions}

The stress state within the Earth's crust is the fundamental driver of earthquakes, yet it cannot be measured directly at depth. Surface geodetic observations---principally GPS velocities and strain-rate azimuths---provide constraints on the \textit{orientation} and \textit{deviatoric pattern} of the stress tensor, but they carry a critical limitation: they cannot constrain the \textit{absolute magnitude} of stress \cite{Khorrami2019}. A region under 10~MPa or 100~MPa of differential stress produces identical surface strain rates.

This \textbf{magnitude ambiguity} has been recognized as a fundamental barrier in tectonophysics \cite{Heidbach2018}. Traditional finite-element approaches \cite{Zamani2017, Rajabi2017} attempt to overcome it by manually prescribing boundary conditions---plate-edge tractions or displacements---but these are themselves poorly constrained and introduce subjective bias.

Recently, Masson et al.\ \cite{Masson2023} demonstrated that a physics-informed neural network can bypass explicit boundary conditions entirely, recovering the 2D elastic stress field of Australia from GPS velocities and World Stress Map orientations. Their work represents a significant advance, but remains limited to: (i) 2D plane stress, (ii) homogeneous effective elasticity, and (iii) purely kinematic constraints with no link to seismic hazard.

We propose to resolve the magnitude ambiguity by introducing a \textbf{third data source}---the historical earthquake catalog---coupled to the stress field through established rate-and-state friction physics.

\section{Governing Equations}

\subsection{Conservation of Momentum ($\mathcal{L}_{\text{PDE}}$)}
The stress tensor $\sigma_{ij}$ must satisfy static equilibrium:
\begin{equation}
    \frac{\partial \sigma_{ij}}{\partial x_j} + \rho(\mathbf{x})\, g_i = 0
    \label{eq:momentum}
\end{equation}
where $\rho(\mathbf{x})$ is the spatially varying crustal density derived from 3D P-wave tomography via the empirical relation $\rho = 2500 + 200 \cdot (V_p - 5)$ \cite{BroNoy2008}, and $g_i$ is the gravitational acceleration vector.

\subsection{Constitutive Law ($\mathcal{L}_{\text{const}}$)}
We adopt a \textbf{steady-state viscous} (Stokes) constitutive relationship for the interseismic period:
\begin{equation}
    \tau_{ij} = 2\,\eta(\mathbf{x})\,\dot{\varepsilon}^{\prime}_{ij}
    \label{eq:constitutive}
\end{equation}
where $\tau_{ij}$ is the deviatoric stress tensor, $\dot{\varepsilon}^{\prime}_{ij}$ is the deviatoric strain-rate tensor, and $\eta(\mathbf{x})$ is the spatially varying effective viscosity. This is the long-timescale limit of Maxwell viscoelasticity (i.e., $\dot{\sigma}_{ij} \to 0$ for $t \gg \tau_M = \eta/\mu$). For the Iranian lithosphere with $\mu \sim 30$~GPa and $\eta \sim 10^{21}$~Pa$\cdot$s, the Maxwell relaxation time $\tau_M \sim 10^{3}$~yr. Since we model the quasi-static interseismic stress field averaged over timescales of $\sim 10^{4}$--$10^{5}$~yr (consistent with the tectonic strain rates of $\sim 10^{-15}$~s$^{-1}$ from GPS \cite{Khorrami2019}), the Stokes approximation is formally justified.

The effective viscosity $\eta(\mathbf{x})$ is derived from the tomographic shear modulus $\mu(\mathbf{x})$ via an assumed Maxwell relaxation time $\tau_M(\mathbf{x})$:
\begin{equation}
    \eta(\mathbf{x}) = \mu(\mathbf{x}) \cdot \tau_M(\mathbf{x})
    \label{eq:viscosity}
\end{equation}
In the initial implementation, $\tau_M$ is treated as a depth-dependent parameter. A uniform $\tau_M$ is a known oversimplification; future work will incorporate temperature-dependent rheology.

\subsection{Coordinate System and Normalization}
\label{sec:coords}
All spatial coordinates are normalized to the half-open cube $[-1, 1]^2 \times [-1, 0]$. The $x$ and $y$ axes are isotropic in UTM meters; the $z$ axis uses the convention $z = 0$ at the Earth's surface and $z = -1$ at the maximum depth of the velocity model. This $z$-asymmetry is deliberate: the network is evaluated only on the physical crustal volume $\{z \in [-1, 0]\}$, the traction-free surface boundary condition applies at $z = 0$, and the gravitational body force acts in the $-z$ direction. Using a symmetric $[-1, 1]$ range in $z$ would force the network to represent the unphysical half-space $z > 0$ (above the surface) and would double the collocation domain without adding information.

The bounding box is taken as the union of the GPS, declustered catalog, and tomography coverage (after projecting GPS and catalog latitudes/longitudes to the same UTM zone as the velocity model). The characteristic length scale $L_0 = 10^6$~m is half the maximum horizontal extent of this bounding box, and the characteristic depth scale $Z_0$ is the maximum depth of the velocity model. Surface collocation points for the traction-free boundary condition are placed at $z = 0$; volume collocation points are sampled uniformly on $[-1, 0]$ in $z$.

\subsection{Kinematic Data Constraints ($\mathcal{L}_{\text{data}}$)}
GPS strain-rate azimuths $\theta_{\text{obs}}$ constrain the predicted maximum horizontal stress orientation $\theta_{\text{pred}}$ at the surface. The azimuthal loss handles the $180^\circ$ ambiguity of the stress tensor \cite{Heidbach2018}:
\begin{equation}
    \mathcal{L}_{\text{data}} = \frac{1}{N_{\text{GPS}}} \sum_{k=1}^{N_{\text{GPS}}} \sin^2\!\left(2\left(\theta_{\text{pred}}^{(k)} - \theta_{\text{obs}}^{(k)}\right)\right)
    \label{eq:data_loss}
\end{equation}
where $\theta_{\text{pred}} = \frac{1}{2}\arctan\!\left(\frac{2\sigma_{xy}}{\sigma_{xx} - \sigma_{yy}}\right)$.

\subsection{Coulomb--Dieterich Seismicity Coupling ($\mathcal{L}_{\text{seis}}$)}

This is the central theoretical contribution. Rather than treating seismicity as a statistical phenomenon disconnected from the stress field, we embed the established rate-and-state friction framework of Dieterich \cite{Dieterich1994} directly into the PINN loss.

\subsubsection{Coulomb Failure Function}
Given the full stress tensor $\sigma_{ij}(\mathbf{x})$ predicted by the PINN, we compute the Coulomb Failure Function (CFF) on optimally oriented fault planes:
\begin{equation}
    \text{CFF}(\mathbf{x}) = \tau_{\max}(\mathbf{x}) - \mu_f \left(\sigma_n(\mathbf{x}) - P_f\right)
    \label{eq:cff}
\end{equation}
where:
\begin{itemize}
    \item $\tau_{\max}$ is the maximum resolved shear stress on optimally oriented planes, computed from the eigenvalues of $\sigma_{ij}$: $\tau_{\max} = (\sigma_1 - \sigma_3)/2$
    \item $\sigma_n = (\sigma_1 + \sigma_3)/2$ is the mean normal stress on those planes
    \item $\mu_f = 0.6$ is the friction coefficient \cite{Byerlee1978}
    \item $P_f$ is the pore fluid pressure, parameterized as a fraction of lithostatic: $P_f = \lambda_p \cdot \rho g z$ with $\lambda_p \in [0.37, 0.7]$
\end{itemize}

\subsubsection{Rate-and-State Seismicity Rate}
The Dieterich \cite{Dieterich1994} model predicts the seismicity rate as a function of the Coulomb stress state:
\begin{equation}
    R(\mathbf{x}) = r_0 \cdot \exp\!\left(\frac{\text{CFF}(\mathbf{x})}{A\,\sigma_n}\right)
    \label{eq:dieterich}
\end{equation}
where $r_0$ is the reference background seismicity rate and $A \sim 0.005$--$0.01$ is the rate-and-state constitutive parameter. Note that this formulation uses the \textit{Coulomb failure function} on optimally oriented planes---not the raw differential stress---which correctly accounts for the Mohr--Coulomb failure criterion.

\subsubsection{Poisson Log-Likelihood Loss}
The earthquake catalog is a marked spatio-temporal point process. The natural loss function is the negative Poisson log-likelihood \cite{Ogata1998}:
\begin{equation}
    \mathcal{L}_{\text{seis}} = -\sum_{i=1}^{N_{\text{eq}}} \log R(\mathbf{x}_i) + \int_\Omega R(\mathbf{x})\,d\mathbf{x}
    \label{eq:poisson}
\end{equation}
Both terms are evaluated in practice by Monte Carlo. The first term is averaged over the $N_{\text{eq}}$ declustered catalog events with $M \geq M_c$, where $M_c$ is the magnitude of completeness estimated via the maximum curvature method \cite{Wiemer2000}. The second term (the integral) is approximated by uniform sampling of the domain volume:
\begin{equation}
    \mathcal{L}_{\text{seis}} \;\approx\; -\frac{1}{N_{\text{eq}}}\sum_{i=1}^{N_{\text{eq}}} \log R(\mathbf{x}_i) \;+\; \frac{V_\Omega}{N_{\text{coll}}} \sum_{j=1}^{N_{\text{coll}}} R(\mathbf{x}_j)
    \label{eq:poisson_mc}
\end{equation}
where $V_\Omega = (X_{\max} - X_{\min})(Y_{\max} - Y_{\min})(Z_{\max} - Z_{\min})$ is the physical domain volume in m$^3$, computed from the union of the GPS, catalog, and tomography bounding boxes (after coordinate projection to a common reference frame), and $N_{\text{coll}}$ is the number of physics collocation points sampled per training step. The $1/N_{\text{eq}}$ normalization on the first term ensures that the data-fit and integral contributions are on comparable scales during loss-weight tuning; without it, the magnitude of the first term would scale with catalog size and dominate $w_{\text{seis}}$ selection.

\subsection{Total Loss}
\begin{equation}
    \mathcal{L}_{\text{total}} = w_{\text{pde}}\,\mathcal{L}_{\text{PDE}} + w_{\text{const}}\,\mathcal{L}_{\text{const}} + w_{\text{data}}\,\mathcal{L}_{\text{data}} + w_{\text{bc}}\,\mathcal{L}_{\text{BC}} + w_{\text{seis}}\,\mathcal{L}_{\text{seis}}
    \label{eq:total_loss}
\end{equation}
where $\mathcal{L}_{\text{BC}}$ enforces the traction-free surface boundary condition ($\sigma_{xz} = \sigma_{yz} = \sigma_{zz} = 0$ at $z = 0$).

\section{Network Architecture}

\begin{table}[h!]
\centering
\caption{Comparison with Masson et al.\ \cite{Masson2023} (ML-SEISMIC).}
\begin{tabular}{lcc}
\toprule
\textbf{Feature} & \textbf{ML-SEISMIC} & \textbf{This Work} \\
\midrule
Dimensionality & 2D & \textbf{3D} \\
Constitutive Law & Linear Elastic & \textbf{Viscous (Stokes)} \\
Material Properties & Homogeneous (trained) & \textbf{Heterogeneous (tomography)} \\
Input Embedding & None & \textbf{Fourier Features} \\
Seismicity Coupling & None & \textbf{Coulomb--Dieterich} \\
Boundary Conditions & 2 displacement pins & \textbf{Traction-free surface} \\
Validation & Analytical benchmark & \textbf{Synthetic + WSM + FEM} \\
\bottomrule
\end{tabular}
\label{tab:comparison}
\end{table}

The neural network $f_\theta : \mathbb{R}^3 \to \mathbb{R}^9$ maps normalized spatial coordinates $(x, y, z)$ to the velocity and stress fields $(v_x, v_y, v_z, \sigma_{xx}, \sigma_{yy}, \sigma_{zz}, \sigma_{xy}, \sigma_{yz}, \sigma_{xz})$. The architecture employs:
\begin{itemize}
    \item \textbf{Fourier Feature Mapping} \cite{Tancik2020}: random Fourier projection $\gamma(\mathbf{x}) = [\sin(2\pi\mathbf{B}\mathbf{x}), \cos(2\pi\mathbf{B}\mathbf{x})]$ to overcome the spectral bias of standard MLPs, enabling resolution of high-frequency stress gradients near fault zones.
    \item \textbf{Fully-connected MLP}: 4 hidden layers with $\tanh$ activation and Xavier initialization.
    \item \textbf{Non-dimensionalization}: all quantities are scaled by characteristic values ($S_0 = 1$~GPa, $V_0 = 10^{-9}$~m/s, $L_0 = 10^6$~m) to ensure $O(1)$ loss residuals. Spatial coordinates follow the asymmetric convention $[-1, 1]^2 \times [-1, 0]$ described in Section~\ref{sec:coords}.
\end{itemize}

\section{Data Sources}

\begin{enumerate}
    \item \textbf{GPS strain-rate azimuths}: $\sim$2000 observations from Khorrami et al.\ \cite{Khorrami2019} and Rayisi \cite{Rayisi2016}, providing maximum compressive strain-rate orientation across the Iranian Plateau.
    \item \textbf{3D P-wave velocity model}: Tomographic model covering Central Iran to $\sim$30~km depth, providing spatially heterogeneous $V_p(\mathbf{x})$, from which $\rho$, $\mu$, and $\eta$ are derived.
    \item \textbf{Earthquake catalog}: Historical and instrumental seismicity of Central Iran ($M_w \geq M_c$), used to construct the spatial seismicity rate for the coupling loss. Aftershock declustering is applied via the Zaliapin--Ben-Zion method \cite{Zaliapin2008}.
\end{enumerate}

\section{Validation Plan}

\subsection{Synthetic Recovery Test}
A known 3D stress field is generated from an analytical or FEM solution. Synthetic GPS azimuths and a synthetic earthquake catalog (sampled from the Coulomb--Dieterich rate) are produced. The PINN must recover the input stress field to within a defined tolerance. This demonstrates well-posedness: the combination of surface azimuths and volumetric seismicity data provides sufficient constraint for a unique solution.

\subsection{Independent Validation}
\begin{enumerate}
    \item \textbf{World Stress Map}: predicted $S_{H\text{max}}$ orientations are compared against independent WSM records for Iran \cite{Heidbach2018}.
    \item \textbf{FEM comparison}: results are compared against the finite-element stress model of Zamani \cite{Zamani2017} for Central Iran.
    \item \textbf{Leave-one-out cross-validation}: GPS stations are held out individually; the model must predict their azimuth from the remaining data.
\end{enumerate}

\subsection{Uncertainty Quantification}
An ensemble of PINNs with different random initializations and Fourier frequency matrices provides epistemic uncertainty estimates on predicted stress magnitudes and seismicity rates.

\section{Scientific Significance}

The coupling term $\mathcal{L}_{\text{seis}}$ (Eq.~\ref{eq:poisson}) resolves the magnitude ambiguity through a mechanism grounded in three decades of earthquake physics \cite{Dieterich1994}:
\begin{itemize}
    \item Regions with frequent seismicity anchor the Coulomb stress level to values near the failure threshold, constraining the \textit{absolute magnitude} of the stress tensor.
    \item Aseismic regions are constrained to have Coulomb stress below the failure threshold, or to lack suitably oriented faults.
    \item The rate-and-state framework naturally predicts the Omori aftershock decay law, providing internal consistency checks.
\end{itemize}

No previous work has coupled a PINN-based stress inversion with rate-and-state seismicity physics. This represents a unification of continuum mechanics (tectonophysics) with statistical seismology within a single differentiable framework.

\section{Expected Impact}

This framework produces the first continuous, physics-constrained 3D map of both stress tensor and seismogenic potential for an active collision zone. Unlike statistical hazard maps that extrapolate from past seismicity patterns, this model explains \textit{why} earthquakes occur where they do by computing the Coulomb failure state from first principles, constrained by surface geodesy and deep tomographic structure.

\begin{thebibliography}{99}

\bibitem{Masson2023}
Masson, Y., Lhomme, T. \& Tillier, E. (2023). Physics-informed neural network reconciles Australian displacements and tectonic stresses. \textit{Sci.\ Rep.}\ \textbf{13}, 22770.

\bibitem{Dieterich1994}
Dieterich, J.H. (1994). A constitutive law for rate of earthquake production and its application to earthquake clustering. \textit{J.\ Geophys.\ Res.}\ \textbf{99}, 2601--2618.

\bibitem{Byerlee1978}
Byerlee, J.D. (1978). Friction of rocks. \textit{Pure Appl.\ Geophys.}\ \textbf{116}, 615--626.

\bibitem{Ogata1998}
Ogata, Y. (1998). Space-time point-process models for earthquake occurrences. \textit{Ann.\ Inst.\ Stat.\ Math.}\ \textbf{50}, 379--402.

\bibitem{Wiemer2000}
Wiemer, S. \& Wyss, M. (2000). Minimum magnitude of completeness in earthquake catalogs. \textit{Bull.\ Seismol.\ Soc.\ Am.}\ \textbf{90}, 859--869.

\bibitem{Zaliapin2008}
Zaliapin, I. \& Ben-Zion, Y. (2008). Earthquake clusters in southern California. \textit{J.\ Geophys.\ Res.}\ \textbf{113}, B12301.

\bibitem{Tancik2020}
Tancik, M., et al. (2020). Fourier Features Let Networks Learn High Frequency Functions in Low Dimensional Domains. \textit{NeurIPS}.

\bibitem{Karniadakis2021}
Karniadakis, G.E., et al. (2021). Physics-informed machine learning. \textit{Nat.\ Rev.\ Phys.}\ \textbf{3}, 422--440.

\bibitem{Haghighat2021}
Haghighat, E., Raissi, M., Moure, A., Gomez, H. \& Juanes, R. (2021). A physics-informed deep learning framework for inversion and surrogate modeling in solid mechanics. \textit{Comput.\ Methods Appl.\ Mech.\ Eng.}\ \textbf{379}, 113741.

\bibitem{Zhang2022}
Zhang, H., et al. (2022). A physics-informed neural network for modeling fracture without gradient damage. \textit{Sci.\ Rep.}

\bibitem{Li2023}
Li, Y., et al. (2023). Inferring viscoplastic models from velocity fields: a physics-informed neural network approach. \textit{arXiv:2506.17681}.

\bibitem{Degen2023}
Degen, D., et al. (2023). Perspectives of Physics-Based Machine Learning Strategies for Geoscientific Applications. \textit{GMD}.

\bibitem{Zamani2017}
Zamani, G. (2017). Finite Element Modeling of Tectonic Stress Field of Central Iran. \textit{Acta Geodyn.\ Geomater.}

\bibitem{Rajabi2017}
Rajabi, M., Heidbach, O., Tingay, M. \& Reiter, K. (2017). Prediction of the present-day stress field in the Australian continental crust using 3D geomechanical-numerical models. \textit{Aust.\ J.\ Earth Sci.}\ \textbf{64}, 435--454.

\bibitem{Ranalli1995}
Ranalli, G. (1995). \textit{Rheology of the Earth}. Chapman \& Hall.

\bibitem{Khorrami2019}
Khorrami, F., et al. (2019). Kinematic analysis of the Iranian Plateau using GPS data. \textit{GJI}.

\bibitem{Rayisi2016}
Rayisi, A. (2016). GPS strain rate analysis of the Iranian Plateau. \textit{M.Sc.\ Thesis}.

\bibitem{Anderson1951}
Anderson, E.M. (1951). \textit{The Dynamics of Faulting}. Oliver \& Boyd.

\bibitem{Heidbach2018}
Heidbach, O., et al. (2018). The world stress map database release 2016. \textit{Tectonophysics}\ \textbf{744}, 484--498.

\bibitem{BroNoy2008}
Brocher, T.M. (2005). Empirical relations between elastic wavespeeds and density in the Earth's crust. \textit{Bull.\ Seismol.\ Soc.\ Am.}\ \textbf{95}, 2081--2092.

\bibitem{GCMT}
Ekstr\"om, G., et al. (2012). The global CMT project 2004--2010. \textit{PEPI}.

\end{thebibliography}

\end{document}
